% 【拆分生成】 questions.tex —— 2022 年·新高考Ⅰ卷 · 题目内容片段（无答案）
% 【拆分生成】 由 tools/split_project.py 自合并书 papers/2022_xgk1.tex 包裹 question 环境生成，勿手改
% 由 tools/md2tex.py 自 2022_数学真题_试题.md 生成，勿手改（改源 MD 后重跑）


\begin{gknotice}
\textbf{结构提示}：本卷为 \textbf{22 题制}。2024 年起新高考 I 卷改为 19 题制，故本卷宜按知识点抽题使用，不宜整套限时模拟。

\textbf{适用省份}：山东、广东、湖南、湖北、河北、江苏、福建
\end{gknotice}

\par\medskip

\gkpart{一、单选题}

\begin{gktip}
本题共 8 小题，每小题 5 分，共 40 分。在每小题给出的四个选项中，只有一项是符合题目要求的。
\end{gktip}

\begin{question}
\gkti{1}若集合 $M=\left\{x\mid\sqrt x<4\right\}$，$N=\{x\mid3x\ge1\}$，则 $M\cap N=$（　　）

\begin{gkopts}(2)
\task A．$\{x\mid0\le x<2\}$
\task B．$\left\{x\ \middle|\ \dfrac13\le x<2\right\}$
\task C．$\{x\mid3\le x<16\}$
\task D．$\left\{x\ \middle|\ \dfrac13\le x<16\right\}$
\end{gkopts}

\end{question}
\begin{question}
\gkti{2}若 $\mathrm i(1-z)=1$，则 $z+\bar z=$（　　）

\begin{gkopts}(4)
\task A．$-2$
\task B．$-1$
\task C．$1$
\task D．$2$
\end{gkopts}

\end{question}
\begin{question}
\gkti{3}在 $\triangle ABC$ 中，点 $D$ 在边 $AB$ 上，$BD=2DA$。记 $\overrightarrow{CA}=\boldsymbol m$，$\overrightarrow{CD}=\boldsymbol n$，则 $\overrightarrow{CB}=$（　　）

\begin{gkopts}(2)
\task A．$3\boldsymbol m-2\boldsymbol n$
\task B．$-2\boldsymbol m+3\boldsymbol n$
\task C．$3\boldsymbol m+2\boldsymbol n$
\task D．$2\boldsymbol m+3\boldsymbol n$
\end{gkopts}

\end{question}
\begin{question}
\gkti{4}南水北调工程缓解了北方一些地区水资源短缺问题，其中一部分水蓄入某水库。已知该水库水位为海拔 $148.5\ \mathrm m$ 时，相应水面的面积为 $140.0\ \mathrm{km}^{2}$；水位为海拔 $157.5\ \mathrm m$ 时，相应水面的面积为 $180.0\ \mathrm{km}^{2}$。将该水库在这两个水位间的形状看作一个棱台，则该水库水位从海拔 $148.5\ \mathrm m$ 上升到 $157.5\ \mathrm m$ 时，增加的水量约为（$\sqrt7\approx2.65$）（　　）

\begin{gkopts}(2)
\task A．$1.0\times10^{9}\ \mathrm m^{3}$
\task B．$1.2\times10^{9}\ \mathrm m^{3}$
\task C．$1.4\times10^{9}\ \mathrm m^{3}$
\task D．$1.6\times10^{9}\ \mathrm m^{3}$
\end{gkopts}

\end{question}
\begin{question}
\gkti{5}从 $2$ 到 $8$ 的 $7$ 个整数中随机取 $2$ 个不同的数，则这 $2$ 个数互质的概率为（　　）

\begin{gkopts}(4)
\task A．$\dfrac16$
\task B．$\dfrac13$
\task C．$\dfrac12$
\task D．$\dfrac23$
\end{gkopts}

\end{question}
\begin{question}
\gkti{6}记函数 $f(x)=\sin\left(\omega x+\dfrac\pi4\right)+b\ (\omega>0)$ 的最小正周期为 $T$。若 $\dfrac{2\pi}3<T<\pi$，且 $y=f(x)$ 的图象关于点 $\left(\dfrac{3\pi}2,2\right)$ 中心对称，则 $f\left(\dfrac\pi2\right)=$（　　）

\begin{gkopts}(4)
\task A．$1$
\task B．$\dfrac32$
\task C．$\dfrac52$
\task D．$3$
\end{gkopts}

\end{question}
\begin{question}
\gkti{7}设 $a=0.1\mathrm e^{0.1}$，$b=\dfrac19$，$c=-\ln0.9$，则（　　）

\begin{gkopts}(4)
\task A．$a<b<c$
\task B．$c<b<a$
\task C．$c<a<b$
\task D．$a<c<b$
\end{gkopts}

\end{question}
\begin{question}
\gkti{8}已知正四棱锥的侧棱长为 $l$，其各顶点都在同一球面上。若该球的体积为 $36\pi$，且 $3\le l\le3\sqrt3$，则该正四棱锥体积的取值范围是（　　）

\begin{gkopts}(2)
\task A．$\left[18,\dfrac{81}4\right]$
\task B．$\left[\dfrac{27}4,\dfrac{81}4\right]$
\task C．$\left[\dfrac{27}4,\dfrac{64}3\right]$
\task D．$[18,27]$
\end{gkopts}

\par\medskip

\end{question}
\gkpart{二、多选题}

\begin{gktip}
本题共 4 小题，每小题 5 分，共 20 分。全部选对得 5 分，部分选对得 2 分，有选错得 0 分。
\end{gktip}

\begin{question}
\gkti{9}已知正方体 $ABCD\text{-}A_1B_1C_1D_1$，则（　　）

\begin{gkopts}(2)
\task A．直线 $BC_1$ 与 $DA_1$ 所成的角为 $90^\circ$
\task B．直线 $BC_1$ 与 $CA_1$ 所成的角为 $90^\circ$
\task C．直线 $BC_1$ 与平面 $BB_1D_1D$ 所成的角为 $45^\circ$
\task D．直线 $BC_1$ 与平面 $ABCD$ 所成的角为 $45^\circ$
\end{gkopts}

\end{question}
\begin{question}
\gkti{10}已知函数 $f(x)=x^{3}-x+1$，则（　　）

\begin{gkopts}(2)
\task A．$f(x)$ 有两个极值点
\task B．$f(x)$ 有三个零点
\task C．点 $(0,1)$ 是曲线 $y=f(x)$ 的对称中心
\task D．直线 $y=2x$ 是曲线 $y=f(x)$ 的切线
\end{gkopts}

\end{question}
\begin{question}
\gkti{11}已知 $O$ 为坐标原点，点 $A(1,1)$ 在抛物线 $C:x^{2}=2py\ (p>0)$ 上，过点 $B(0,-1)$ 的直线交 $C$ 于 $P$，$Q$ 两点，则（　　）

\begin{gkopts}(2)
\task A．$C$ 的准线为 $y=-1$
\task B．直线 $AB$ 与 $C$ 相切
\task C．$|OP|\cdot|OQ|>|OA|^{2}$
\task D．$|BP|\cdot|BQ|>|BA|^{2}$
\end{gkopts}

\end{question}
\begin{question}
\gkti{12}已知函数 $f(x)$ 及其导函数 $f'(x)$ 的定义域均为 $\mathbf R$，记 $g(x)=f'(x)$。若 $f\left(\dfrac32-2x\right)$，$g(2+x)$ 均为偶函数，则（　　）

\begin{gkopts}(2)
\task A．$f(0)=0$
\task B．$g\left(-\dfrac12\right)=0$
\task C．$f(-1)=f(4)$
\task D．$g(-1)=g(2)$
\end{gkopts}

\par\medskip

\end{question}
\gkpart{三、填空题}

\begin{gktip}
本题共 4 小题，每小题 5 分，共 20 分。
\end{gktip}

\begin{question}
\gkti{13}$\left(1-\dfrac yx\right)(x+y)^{8}$ 的展开式中 $x^{2}y^{6}$ 的系数为 \gkblank{}（用数字作答）。

\end{question}
\begin{question}
\gkti{14}写出与圆 $x^{2}+y^{2}=1$ 和 $(x-3)^{2}+(y-4)^{2}=16$ 都相切的一条直线的方程 \gkblank{}．

\end{question}
\begin{question}
\gkti{15}若曲线 $y=(x+a)\mathrm e^{x}$ 有两条过坐标原点的切线，则 $a$ 的取值范围是 \gkblank{}．

\end{question}
\begin{question}
\gkti{16}已知椭圆 $C:\dfrac{x^{2}}{a^{2}}+\dfrac{y^{2}}{b^{2}}=1\ (a>b>0)$，$C$ 的上顶点为 $A$，两个焦点为 $F_1$，$F_2$，离心率为 $\dfrac12$。过 $F_1$ 且垂直于 $AF_2$ 的直线与 $C$ 交于 $D$，$E$ 两点，$|DE|=6$，则 $\triangle ADE$ 的周长是 \gkblank{}．

\par\medskip

\end{question}
\gkpart{四、解答题}

\begin{gktip}
本题共 6 小题，共 70 分。解答应写出文字说明、证明过程或演算步骤。
\end{gktip}

\begin{question}
\gkti{17}\textbf{（10 分）} 记 $S_n$ 为数列 $\left\{a_n\right\}$ 的前 $n$ 项和，已知 $a_1=1$，$\left\{\dfrac{S_n}{a_n}\right\}$ 是公差为 $\dfrac13$ 的等差数列。

  (1) 求 $\left\{a_n\right\}$ 的通项公式；

  (2) 证明：$\dfrac1{a_1}+\dfrac1{a_2}+\cdots+\dfrac1{a_n}<2$。

\end{question}
\begin{question}
\gkti{18}\textbf{（12 分）} 记 $\triangle ABC$ 的内角 $A$，$B$，$C$ 的对边分别为 $a$，$b$，$c$，已知 $\dfrac{\cos A}{1+\sin A}=\dfrac{\sin2B}{1+\cos2B}$。

  (1) 若 $C=\dfrac{2\pi}3$，求 $B$；

  (2) 求 $\dfrac{a^{2}+b^{2}}{c^{2}}$ 的最小值。

\end{question}
\begin{question}
\gkti{19}\textbf{（12 分）} 如图，直三棱柱 $ABC\text{-}A_1B_1C_1$ 的体积为 $4$，$\triangle A_1BC$ 的面积为 $2\sqrt2$。

\begin{center}% 2022 新高考Ⅰ卷 19：直三棱柱 ABC-A1B1C1（D 为 A1C 的中点）
\begin{tikzpicture}[scale=0.95, font=\small]
  \coordinate (A)  at (0,0.5);
  \coordinate (B)  at (1.15,-0.42);
  \coordinate (C)  at (2.5,0.5);
  \coordinate (A1) at (0,2.7);
  \coordinate (B1) at (1.15,1.78);
  \coordinate (C1) at (2.5,2.7);
  \coordinate (D)  at ($(A1)!0.5!(C)$);
  \draw (A) -- (B) -- (C);
  \draw[dashed] (A) -- (C);
  \draw (A) -- (A1)  (B) -- (B1)  (C) -- (C1);
  \draw (A1) -- (B1) -- (C1) -- cycle;
  \draw (A1) -- (B);
  \draw[dashed] (A1) -- (C);
  \draw[dashed] (B) -- (D);
  \fill (D) circle (0.9pt);
  \node[left] at (A) {$A$};
  \node[below] at (B) {$B$};
  \node[right] at (C) {$C$};
  \node[above left, inner sep=1pt] at (A1) {$A_1$};
  \node[below, inner sep=2.5pt] at (B1) {$B_1$};
  \node[above right, inner sep=1pt] at (C1) {$C_1$};
  \node[above right, inner sep=0.5pt] at (D) {$D$};
\end{tikzpicture}
\end{center}

  (1) 求 $A$ 到平面 $A_1BC$ 的距离；

  (2) 设 $D$ 为 $A_1C$ 的中点，$AA_1=AB$，平面 $A_1BC\perp$ 平面 $ABB_1A_1$，求二面角 $A\text{-}BD\text{-}C$ 的正弦值。

\end{question}
\begin{question}
\gkti{20}\textbf{（12 分）} 一医疗团队为研究某地的一种地方性疾病与当地居民的卫生习惯（卫生习惯分为良好和不够良好两类）的关系，在已患该疾病的病例中随机调查了 $100$ 例（称为病例组），同时在未患该疾病的人群中随机调查了 $100$ 人（称为对照组），得到如下数据：

\begin{center}
\begin{tabular}{lcc}
\toprule
 & 不够良好 & 良好 \\
\midrule
病例组 & $40$ & $60$ \\
对照组 & $10$ & $90$ \\
\bottomrule
\end{tabular}
\end{center}

  (1) 能否有 $99\%$ 的把握认为患该疾病群体与未患该疾病群体的卫生习惯有差异？

  (2) 从该地的人群中任选一人，$A$ 表示事件"选到的人卫生习惯不够良好"，$B$ 表示事件"选到的人患有该疾病"。$\dfrac{P(B\mid A)}{P(\bar B\mid A)}$ 与 $\dfrac{P(B\mid\bar A)}{P(\bar B\mid\bar A)}$ 的比值是卫生习惯不够良好对患该疾病风险程度的一项度量指标，记该指标为 $R$。

  　① 证明：$R=\dfrac{P(A\mid B)}{P(\bar A\mid B)}\cdot\dfrac{P(\bar A\mid\bar B)}{P(A\mid\bar B)}$；

  　② 利用该调查数据，给出 $P(A\mid B)$，$P(A\mid\bar B)$ 的估计值，并利用①的结果给出 $R$ 的估计值。

  附：$K^{2}=\dfrac{n(ad-bc)^{2}}{(a+b)(c+d)(a+c)(b+d)}$

\begin{center}
\begin{tabular}{cccc}
\toprule
$P\left(K^{2}\ge k\right)$ & $0.050$ & $0.010$ & $0.001$ \\
\midrule
$k$ & $3.841$ & $6.635$ & $10.828$ \\
\bottomrule
\end{tabular}
\end{center}

\end{question}
\begin{question}
\gkti{21}\textbf{（12 分）} 已知点 $A(2,1)$ 在双曲线 $C:\dfrac{x^{2}}{a^{2}}-\dfrac{y^{2}}{a^{2}-1}=1\ (a>1)$ 上，直线 $l$ 交 $C$ 于 $P$，$Q$ 两点，直线 $AP$，$AQ$ 的斜率之和为 $0$。

  (1) 求 $l$ 的斜率；

  (2) 若 $\tan\angle PAQ=2\sqrt2$，求 $\triangle PAQ$ 的面积。

\end{question}
\begin{question}
\gkti{22}\textbf{（12 分）} 已知函数 $f(x)=\mathrm e^{x}-ax$ 和 $g(x)=ax-\ln x$ 有相同的最小值。

  (1) 求 $a$；

  (2) 证明：存在直线 $y=b$，其与两条曲线 $y=f(x)$ 和 $y=g(x)$ 共有三个不同的交点，并且从左到右的三个交点的横坐标成等差数列。

\par\medskip


\end{question}